% page 157 of the facsimile (image p-156 of the scan)
% block 162.6 x 250.6 mm ; left margin 39.4 mm
\pgc{17.780mm}{0.000mm}{
\l{}
\l{}
\l{}
\l{}
\l{}
\l{etamino. Stofo ek lano lejera. - Lanajo, krinajo,quan onu uzas}
\l{\cel{3}por blutar. - DEFIRS.}
\l{}
\l{etapar. (netrans.) (\sou{pri armeo}). Haltar en loko ube la trupi}
\l{\cel{3}marchanta repozos dum nokto. - DFIRS.}
\l{}
\l{etato. To quo konstatas la estado di la kozi en ta o ca ins-}
\l{\cel{3}tanto. - DFIS.}
\l{}
\l{\sou{etero}. I. La parto maxim subtila di atmosfero, ye qua (krede}
\l{\cel{3}da la antiqui) konsistas fairo. - II. (\sou{fiziko}).Fluido ne-}
\l{\cel{3}ponderebla, di qua ula fizikisti konjektis la existo por}
\l{\cel{3}explikar la fenomeni \textquotedbl{}foto\textquotedbl{} e \textquotedbl{}kaloro\textquotedbl{}. - III. (\sou{kemi}o).}
\l{\cel{3}Liquido volatila qua produktas la alkoholi kande onu de-}
\l{\cel{3}prenas de li equivalanto ek aquo. - DEFIRS.}
\l{}
\l{\sou{eterna}. I. Qua nek komencis nek finos. - II. Qua ne finos. -}
\l{\cel{3}III. Di qua la fino-instanto ne esas previdebla. - EFIS.}
\l{}
\l{\sou{etiko}. I. La cienco pri la mori. - II. Traktato pri ta cienco.}
\l{\cel{3}- DEFIRS.}
\l{}
\l{\sou{etiketo}. Indiko-texto fixigita sur objekto, por savigar la}
\l{\cel{3}naturo, la kontenajo, la preco, la destineso, e c. di ica.}
\l{}
\l{}
\l{\sou{etimologio}. Decendo de vorto relate un o plura vorti altra de}
\l{\cel{3}qui lu derivesas. - DEFIRS.}
\l{}
\l{\sou{etiolar}. (\sou{trans}.) I. Igar (planto) febla e senkolora per kresk-}
\l{\cel{3}igar lu en loko obskura. - II. Igar (persono) febla e pala,}
\l{\cel{3}per privacar lu de aero libera, de korp-exercado. - EFL.}
\l{}
\l{\sou{etiologio}. I. (\sou{filoz.}) Exploro di la kauzi. - II. (\sou{medic}.)}
\l{\cel{3}La parto di medicino, en qua onu exploras la morbi. - DEFIRS.}
\l{}
\l{\sou{etmoido}. (\sou{anat}.) (Osto kraniala) di qua la lamo supra esas}
\l{\cel{3}truoza, inkastrita en la sesguro infra di la osto frontala,}
\l{\cel{3}e qua kontributas a formacar la naztrui. - DEF.}
\l{}
\l{\sou{etnografio}. Deskripto di la populi. - DEFIRS.}
\l{}
\l{\sou{etnologio}. Cienco di la origino e di la deveno di la populi.-}
\l{\cel{3}DEFIRS.}
\l{}
\l{\sou{etologio}. Cienco di la mori. - Traktato pri la mori.}
\l{}
\l{\sou{etuyo}. Sorto di buxo adaptita a la formo di la objekto quan}
\l{\cel{3}lu kontenos. - DFIS.}
\l{}
\l{\sou{eudiometro}. Instrumento uzata por determinar la proporcion}
\l{\cel{3}di la parti qui kompozas aero od irgaltra mixuro de gasi. -}
\l{\cel{3}DEFIRS.}
}
